\section{Experiments: the Main Construction in the model}\label{sec:experiments}

Code: \texttt{experiments/model.py}, \texttt{mc\_flow.py}, \texttt{mc\_profile.py}.
See \texttt{experiments/README.md} for the environment and run commands.
These are numerical observations, not proofs.

\subsection*{Setup}
$\kappa=\frac12$; $u_0$ constant; $u_1$ a normalised bump on $[0,w]$, $w=5\cdot10^{-4}$, made
orthogonal to $u_0$, so $\norm{T^*u_1}=0.009$; $x_0$ as in \eqref{eq:x0}. Functions are stored on
440 Gauss--Legendre nodes. The start is that of Lemma~\ref{lem:start} followed by the
re-extremalisation of Proposition~\ref{prop:domination}:
$\tau=1.9\cdot10^{-3}$, $\norm{x_0-v}=0.498$, $\abs{\hat a_0}^2/\tau=0.998$.

\begin{remark}
In the model this start is close to a true solution, namely
$x_0\mathbf 1_{(w',\kappa]}$ with $w'\approx w$: $\dist(v,\Sol)=0.051$. This is not an accident:
$x_0-v\approx\frac12u_1$ and $u_1$ is nearly in $\ker T^*$.
\end{remark}

A step replaces $y$ by $(1+r(\lambda))y$ with $r$ a real polynomial of degree $<D$. Each step
asks for $\tau\to(1-2\beta)\tau$ together with the listed constraints, takes the solution of
least norm (in the basis $(\lambda/\kappa)^j$), and corrects it by Newton iteration so that the
targets are met exactly, not only to first order. $\beta$ starts at $0.05$ and is halved when
the Newton iteration fails. ``$s_{\min}$'' is the least singular value of the row-normalised
Jacobian of the constrained quantities, a measure of their independence.

\subsection*{Results after 400 steps}
\begin{center}
\small
\begin{tabular}{@{}p{.30\textwidth}cccccc@{}}
\toprule
Constraints & $D$ & final $\tau$ & $\abs{\hat a_0}^2/\tau$ & $\dist(v,\Sol)$ & $s_{\min}$ & final $\beta$\\
\midrule
(start) & & $1.9\cdot10^{-3}$ & 0.998 & 0.051 & & \\
$\varepsilon$ fixed & 14 & $8.9\cdot10^{-6}$ & 0.50 & 0.096 & 0.28 & $9\cdot10^{-5}$\\
$\varepsilon$ fixed & 24 & $8.6\cdot10^{-6}$ & 0.50 & 0.097 & 0.12 & $1\cdot10^{-4}$\\
$\varepsilon$ fixed, $\hat a_0$ tracks $\sqrt\tau$ & 14 & $6.6\cdot10^{-4}$ & 0.999 & 0.071 & $1.4\cdot10^{-4}$ & $2\cdot10^{-5}$\\
$\varepsilon^2$ falls by $2\beta\tau$, $\hat a_0$ tracks & 14 & $5.9\cdot10^{-4}$ & 0.999 & 0.070 & $9\cdot10^{-4}$ & $4\cdot10^{-4}$\\
same & 24 & $5.7\cdot10^{-4}$ & 0.999 & 0.070 & $7\cdot10^{-4}$ & $4\cdot10^{-4}$\\
\bottomrule
\end{tabular}
\end{center}
The last two rows use the prescription of \paper{} for its ``dependent case'' (after (44)).

\subsection*{Observations}
\begin{enumerate}
\item In every variant the state moves \emph{away} from $\Sol$ while the defect falls. The
profile $g=v/x_0$, which should tend to $1$ on $(w,\kappa]$, instead develops oscillations: in
the first row it ends near $1.38,\ 1.21,\ 0.98,\ 0.93,\ 1.06,\ 0.98,\ 1.13$ at
$\lambda=6\cdot10^{-4},\ 0.03,\ 0.1,\ 0.2,\ 0.3,\ 0.4,\ 0.49$. The defect is being reduced by
letting $g-g^2$ change sign so that its low moments cancel, which is the real-valued form of the
oscillating states of Section~\ref{sec:model}. By Theorem~\ref{thm:rigidity} this can never
reach defect $0$ without eventually travelling at least the current distance to $\Sol$.
\item With the leading coefficient made to track $\sqrt\tau$, as \paper{} (33) requires, the
constraints become nearly dependent ($s_{\min}\sim10^{-3}$ or less) and progress almost stops
after the defect has fallen by a factor of about $3$.
\item Without that tracking, $\abs{\hat a_0}^2/\tau$ falls to about $\frac12$, so dominance of
the leading coefficient is lost, and the admissible step size still collapses.
\end{enumerate}

\subsection*{Caveats}
The Main Construction is not a single algorithm: \paper{} only asserts that suitable steps
exist. We tested least-norm steps in one basis with degree at most $24$, in double precision,
with finite-difference Jacobians. A cleverer selection rule might behave differently, and in the
model one certainly exists if one is allowed to use knowledge of $\Sol$. What the experiment
shows is that the natural reading of the construction does not head towards a solution even
from a start that is close to one.
